\textbf{Seats}
\begin{enumerate}
    \item 8 mins: Intro slides \par
    \textbf{Motivation}
    \item 20 mins: Finish Example 1.3.1 \par
    \textbf{Check piazza}
    \item 30 mins: Finish Example 1.3.2 \par
    \textbf{Check piazza}
    \item 38 mins: Finish Example 1.3.3 \par
    \textbf{Check piazza}
    \item 45 mins: Finish Example 1.3.4 \par 
    \item Skip improper fractions and return next lecture (or at least the exercises) \par
    \textbf{Check piazza}
    \item 50 mins: Wrapped up
\end{enumerate}
\textbf{Seats}

\vspace{10mm}
\textbf{Example 1.3.1}

Factorise the denominator $(x - 4)/[(x+3)(2x-1)]$,
\begin{align*}
    \frac{x - 4}{(x+3)(2x-1)} &= \frac{A}{x+3} + \frac{B}{2x-1},\\
    \therefore x-4 &= A(2x-1) + B(x+3)
\end{align*}
Set $x=-3 \quad \rightarrow \quad A = 1$\\
Set $x=1/2 \quad \rightarrow \quad B = -1$
\begin{align*}
    \frac{x-4}{2x^2+5x-3} = \frac{1}{x+3} - \frac{1}{2x-1}.
\end{align*}

\vspace{10mm}
\textbf{Example 1.3.2}
\begin{align*}
    \frac{2x^2+8x+18}{x(x+2)(x-3)} &= \frac{A}{x} + \frac{B}{x+2} + \frac{C}{x-3},\\
    \therefore 2x^2+8x+18 &= A(x+2)(x-3) + Bx(x-3) + Cx(x+2)
\end{align*}
Set $x=0 \quad \rightarrow \quad -6A=18 \quad \rightarrow \quad A=-3$\\
Set $x=-2 \quad \rightarrow \quad 10B=10 \quad \rightarrow \quad B=1$\\
Set $x=3 \quad \rightarrow \quad 15C=60 \quad \rightarrow \quad C=4$\\
Therefore
\begin{align*}
    \frac{2x^2+8x+18}{x(x+2)(x-3)} &= -\frac{3}{x} + \frac{1}{x+2} + \frac{4}{x-3}.
\end{align*}


\vspace{10mm}
\textbf{Example 1.3.3}
Factorise the denominator $(x+9)/[x(x^2-6x+9)] = (x+9)/[x(x-3)^2]$,
\begin{align*}
\frac{x+9}{x(x-3)^2} &= \frac{A}{x} + \frac{B}{x-3} + \frac{C}{(x-3)^2}\\
x+9 &= A(x-3)^2 + Bx(x-3) + Cx
\end{align*}
Set $x=0 \quad \rightarrow \quad 9A=9 \quad \rightarrow \quad A=1$\\
Set $x=3 \quad \rightarrow \quad 3C=12 \quad \rightarrow \quad C=4$\\
Compare coefficients of $x^2$: $A+B=0 \quad \rightarrow \quad B=-1$\\
Therefore,
\begin{align*}
\frac{x+9}{x(x-3)^2} &= \frac{1}{x} - \frac{1}{x-3} + \frac{4}{(x-3)^2}
\end{align*}

\vspace{10mm}
\textbf{Example 1.3.4}
\begin{align*}
    \frac{x^2}{(x-1)(x^2+3x-3)} &= \frac{A}{x-1} + \frac{Bx+C}{x^2+3x-3}\\
    \therefore x^2 &= A(x^2+3x-3) + (Bx+C)(x-1)\\
    \therefore x^2 &= (A+B)x^2 + (3A+B+C)x - (C+3A)
\end{align*}
Comparing coefficients of $x^2$: $A+B=1 \quad \rightarrow \quad B=1-A$\\
Comparing coefficients of $x^0$: $C+3A=0 \quad \rightarrow \quad C=-3A$\\
Comparing coefficients of $x^1$: $3A+B+C=0 \quad \rightarrow \quad 3A+1-A-3A=0 \quad \rightarrow \quad A=1$\\ 
So, $B = 0$ and $C=-3$, therefore
\begin{align*}
    \frac{x^2}{(x-1)(x^2+3x-3)} = \frac{1}{x-1} - \frac{3}{x^2+3x-3}.
\end{align*}

\vspace{10mm}
\textbf{Example 1.3.5}
Factorise the denominator $3(x+1)^2/(x^2+x-2) = 3(x+1)^2/[(x - 1)(x + 2)]$, then
\begin{align*}
    \frac{3(x+1)^2}{(x - 1)(x + 2)} &= A + \frac{B}{x-1} + \frac{C}{x+2}\\
    \therefore 3(x+1)^2 = A(x-1)(x+2) + B(x+2) + C(x-1)
\end{align*}
Set $x=1 \quad \rightarrow \quad 3B=12 \quad \rightarrow \quad B=4$\\
Set $x=-2 \quad \rightarrow \quad -3C=3 \quad \rightarrow \quad C=-1$\\
Comparing coefficients of $x^2$: $A=3$\\
Therefore
\begin{align*}
    \frac{3(x+1)^2}{x^2+x-2} = 3 + \frac{4}{x-1} - \frac{1}{x+2}
\end{align*}

\vspace{10mm}
\textbf{Example 1.3.6}
\begin{itemize}
    \item Order of numerator is 6 and denominator is 4. So need a quadratic.
    \item Repeated root so need a $(x-2)$ and $(x-2)^2$ denominator.
    \item Irreducible quadratic so need a linear term in numerator of this.
\end{itemize}
Therefore
\begin{align*}
    Ax^2 + Bx + C + \frac{D}{x-2} + \frac{E}{(x-2)^2} + \frac{Fx+G}{x^2-6x+50}.
\end{align*}